% ============================================================
%  Response to Reviewers
%  Energy-Efficient Deep Learning Without Backpropagation:
%  A Rigorous Hardware-Validated Comparison
%  PPAM 2026
% ============================================================
\documentclass[11pt,a4paper]{article}

\usepackage[T1]{fontenc}
\usepackage[utf8]{inputenc}
\usepackage[a4paper,margin=2.4cm]{geometry}
\usepackage{microtype}
\usepackage{siunitx}
\usepackage{booktabs}
\usepackage{enumitem}
\usepackage{xcolor}
\usepackage[colorlinks=true,allcolors=blue]{hyperref}

\sisetup{detect-weight=true, detect-family=true}
\DeclareSIUnit\mebibyte{MiB}
\emergencystretch=3em

\newcommand{\tmma}{\texttt{torch.cuda.\allowbreak max\_memory\_\allowbreak allocated}}

\newcommand{\FF}{FF}
\newcommand{\MF}{MF}
\newcommand{\BP}{BP}
\newcommand{\CaFo}{CaFo}
\newcommand{\CaFoRand}{CaFo-Rand}
\newcommand{\CaFoDFA}{CaFo-DFA}
\newcommand{\BPDS}{BP-DS}

\newenvironment{comment}%
  {\begin{quote}\itshape\small}%
  {\end{quote}}

\newcommand{\where}[1]{\par\smallskip\noindent\textbf{Where:} #1\par}

\setlist[description]{leftmargin=0pt, labelindent=0pt, style=nextline}

\title{\vspace{-1.5cm}Response to Reviewers\\[2pt]
  \large Energy-Efficient Deep Learning Without Backpropagation:\\
  A Rigorous Hardware-Validated Comparison}
\author{}
\date{}

\begin{document}
\maketitle
\vspace{-1.2cm}

\section*{Summary}

We thank all five reviewers for their time and for a set of comments that
turned out to be far more consequential than a normal camera-ready pass. Three
of them (R2c, R2d, R3c/R5d) pointed at exactly the places where our evidence
was weakest, and acting on them changed our headline result. We would rather
report that than defend the accepted version.

The revision rests on a complete re-measurement campaign rather than on
rewriting. Concretely:

\begin{itemize}[nosep,leftmargin=1.4em]
  \item \textbf{Every experiment was re-run under a harmonised stopping rule.}
        The accepted version used a \emph{different} early-stopping criterion
        for each algorithm --- a defect our own audit found and no reviewer
        raised (disclosure~D2 below). Because training time and energy are
        integrals over training duration, an inconsistent stopping rule does
        not perturb an efficiency comparison, it determines it.
  \item \textbf{The headline efficiency claim does not survive.} Under the
        harmonised rule \MF{} is \num{4.0}$\times$ slower than \BP{} and draws
        \num{3.2}$\times$ the energy on CIFAR-10. The published
        ``\num{41}\% less energy, \num{34}\% faster'' \emph{is} reproducible
        under the submitted protocol --- we reproduced it, at
        \num{42.4}\% [\num{38.9}, \num{44.8}] and \num{33.6}\%
        [\num{30.2}, \num{36.0}] --- but its \emph{sign reverses} once the
        stopping rule is equalised. The abstract, the results and the
        conclusion have all been rewritten accordingly.
  \item \textbf{What does survive is the accuracy and the memory result.}
        \MF{} matches or beats \BP{} in every MLP configuration, by
        \numrange{0.01}{5.04}~pp, and holds up to \num{27}\% less per-process
        device memory.
  \item \textbf{Seed counts rose from three to \numrange{5}{122} per
        contrast}, and every accuracy comparison is now an equivalence test
        against a preregistered $\pm\num{0.25}$~pp margin with bootstrap
        confidence intervals.
  \item \textbf{A five-rung ablation ladder from \BP{} to \MF{}} now
        attributes the cost, including the deep-supervision control that R3
        and R5 asked for.
  \item \textbf{We voluntarily disclose five further defects} found by our own
        audit, none of them raised in review, under \emph{Voluntary
        disclosures} below. The repository is public and cited in the paper, so
        all five are discoverable by any reader after publication.
\end{itemize}

\noindent
All numbers below are reproducible from the public repository. Section,
table and figure references are to the \emph{revised} manuscript, which is
\textbf{15 pages}, the PPAM limit.

\bigskip
\noindent\textbf{Revised structure.}
\S1~Introduction \textbar{} \S2~Background and Related Work \textbar{}
\S3~Framework for Fair and Rigorous Comparison \textbar{}
\S4~Experiments and Results (4.1 \FF{}, 4.2 \CaFo{}, 4.3 \MF{},
4.4 Where the Cost Comes From) \textbar{} \S5~Discussion (5.1 Limitations)
\textbar{} \S6~Conclusion.
Table~1 = stopping rules; Table~2 = headline summary for all three algorithms;
Figure~1 = \MF{} hardware trace; Figure~2 = ablation-ladder waterfall;
Figure~3 = equivalence forest.

% ============================================================
\section*{Reviewer 1}
% ============================================================

\subsection*{R1a --- ``Reads like a very condensed technical report \dots\
detrimental to readability for the average PPAM reader.''}

Accepted, and this drove a structural rewrite rather than a cosmetic one. The
former Sections~2 and 3 were merged into a single \S2, Background and Related
Work, so that each algorithm is introduced once, in prose, alongside the work
it descends from. The contributions list became a paragraph. \S4 now leads
with the result and follows with the evidence rather than the reverse.

The discussion, which R2 also found thin, was rebuilt around three named
questions --- why \MF{} is accurate, which memory instrument one should
believe, and what the traces imply for hardware co-design --- instead of a
list of observations.

We also cut hard, because the revision adds a stopping-rule table, an
ablation-ladder subsection and an equivalence analysis to a paper that already
exceeded the limit. The figure count went from eight to three, the Optuna
search-space table and the full environment listing moved to the public
repository (cited by URL from the Code and data paragraph), and the three
figures we kept were chosen because each carries an argument the prose cannot.

\where{\S2 (merged); \S4.4; \S5; the manuscript is now \num{15} pages.}

\subsection*{R1b --- ``The plots in Fig.~3 and Fig.~4 are completely
illegible.''}

Accepted without reservation. Both figures are gone. They were screenshots of
a Weights~\&~Biases dashboard, exported at browser resolution, with the
dashboard's own font sizes --- unreadable at LNCS column width, and, as it
turned out, not reliable either (disclosure~D5).

Every figure in the revision is generated from the measurement archive by a
committed script, rendered as vector PDF at exactly the LNCS text width
(\num{345.827}~pt) so that no scaling is applied at inclusion time, with a
minimum font size of \num{8}~pt and a colour-blind-safe palette that remains
legible in greyscale. Figure~3, the equivalence forest, replaces the old
Figure~4 and carries \num{28} contrasts with their confidence intervals in the
space the screenshot used for one panel.

\where{Figures~1--3; \texttt{scripts/make\_figures.py} in the repository.}

% ============================================================
\section*{Reviewer 2}
% ============================================================

\subsection*{R2a --- ``Lack of details and analysis in the main article
(specifically Section~6) made the paper seem preliminary.''}

Accepted. The discussion was thin because we had description and no
attribution. It now has both.

\S4.4, \emph{Where the Cost Comes From}, decomposes the \BP{}$\to$\MF{}
transition into five single-change rungs: detach the gradient, attach
auxiliary supervision, replace the readout, cache activations on device, cache
them on host. Measured across all four configurations, \textbf{gradient
locality alone accounts for \numrange{228}{272}~pp of \BP{}'s energy};
auxiliary supervision and the readout together move the total by at most
\num{40}~pp and partly cancel on CIFAR; caching contributes nothing. The
efficiency result is therefore attributed to a specific architectural
property, not to an algorithm taken as a whole.

\S5 then explains the mechanism. \textbf{No configuration we measured exceeds
${\approx}\SI{38}{\percent}$ mean GPU utilisation}, and \MF{} reaches only
\SI{7.3}{\percent} against \BP{}'s \SI{16.2}{\percent} on CIFAR-10. These
workloads are bound by kernel launch latency, not by arithmetic: the
layer-sequential schedule issues many small kernels and waits, so the device
idles while the run stays resident and the fixed power floor is integrated
over four times the duration. That single observation explains the energy
result, the absence of any gain from activation caching, and why the
conclusion is platform-conditional.

\where{\S4.4; Figure~2; \S5, all three paragraphs.}

\subsection*{R2b --- ``The details of the different ML architectures are
important \dots\ these details are not discussed much in the paper.''}

Accepted. \S3 now carries a dedicated paragraph, \emph{Architecture, optimiser
and implementation}, giving: initialisation (Kaiming uniform for all weights
and projection matrices); optimisers (AdamW for \BP{} and \BPDS{}, Adam on a
local cross-entropy loss for \CaFo{} predictors and \MF{} layers); MLP
topology (fully connected, ReLU, \numrange{2}{4} hidden layers of \num{1000}
or \num{2000} units); the exact \CaFo{} block ($3{\times}3$ convolution, ReLU,
$2{\times}2$ max-pooling, batch normalisation, \num{32}/\num{128}/\num{512}
channels, three block predictors summed); and the \FF{} goodness threshold
$\theta=2.0$ with peer normalisation (factor \num{0.03}, momentum \num{0.9}).

The paragraph also states what we think is the load-bearing implementation
fact for a hardware paper: every algorithm is built on the same
\texttt{nn.Linear} and \texttt{nn.Conv2d} primitives in one codebase and
differs only in the loss construction and in where \texttt{detach()} is
called, so no method benefits from a fused or hand-written kernel the others
lack. Data augmentation is now stated explicitly and is identical between each
algorithm and its baseline.

\where{\S3, \emph{Architecture, optimiser and implementation}; \S3,
\emph{Experimental Setup}.}

\subsection*{R2c --- ``The paper claimed that the BP baselines were identical
but algorithm-specific components were removed. The removal of these
components may be the reason why certain architectures were better one way or
another.''}

\textbf{This is the most valuable comment we received, and it is correct.} We
did not argue with it; we measured it.

The ablation ladder in \S4.4 exists precisely to answer it. It puts the
removed components back, one at a time, on top of \BP{}, and measures accuracy
and energy at each rung. The result is that the removal is \emph{not} what
drives the differences: auxiliary supervision and the projection readout
together account for at most \num{40}~pp of the energy change and partly
cancel on CIFAR, whereas gradient detachment accounts for
\numrange{228}{272}~pp.

On accuracy the answer is a split, and we report it as such rather than
picking the convenient half. \MF{} minus \BPDS{}:

\begin{center}\small
\begin{tabular}{lrrrl}
\toprule
Configuration & \MF{}~$-$~\BPDS{} & 95\% CI & $n$ & Verdict \\
\midrule
MNIST 2$\times$1000          & $+0.11$~pp & $[-0.01, +0.23]$ & 20 & equivalent \\
Fashion-MNIST 2$\times$1000  & $-0.06$~pp & $[-0.17, +0.06]$ & 48 & equivalent \\
CIFAR-10 3$\times$2000       & $+5.56$~pp & $[+5.37, +5.74]$ & 20 & different \\
CIFAR-100 3$\times$2000      & $+6.10$~pp & $[+5.88, +6.31]$ & 20 & different \\
\bottomrule
\end{tabular}
\end{center}

On the MNIST-scale tasks, auxiliary supervision alone reproduces \MF{}'s
accuracy and locality buys nothing. On CIFAR it does not, and \MF{}'s margin
is large and fully resolved. A further finding falls out of the same
experiment: \textbf{\BPDS{} is \emph{worse} than plain \BP{} on CIFAR}
($-\num{1.65}$ and $-\num{1.06}$~pp, both resolved), so auxiliary heads are
not a free improvement one can bolt onto \BP{}.

\where{\S4.4, \emph{Does locality buy anything deep supervision does not?};
Figure~2; Figure~3(a); \S3 protocol~2, which now states the objection and
points at the measurement.}

\subsection*{R2d --- ``As backprop algorithms are more widely used, the
functionalities in existing software libraries could have been better
optimised \dots\ than forward-only algorithms.''}

\textbf{The hypothesis was correct in a specific and measurable way, and we
tested it rather than arguing about it.}

\MF{}'s layer-sequential schedule recomputes the activations of every
preceding layer at each stage --- an $O(L^2)$ cost that a mature
implementation would remove by caching. So we implemented the cache, in two
variants (device-resident and host-resident), and verified it against the
recomputing path: over \num{1200} and \num{3600} optimiser steps the maximum
absolute difference in both loss and weights was \textbf{exactly zero}. The
optimisation is bit-exact.

It bought \textbf{no speedup}: \numrange{1.00}{1.02}$\times$ on device,
\num{0.97}$\times$ on host, and the device variant costs roughly \num{40}\%
more per-process memory. The missing optimisation was not missing; it simply
does not pay, for the reason given under R2a --- at these layer widths the
device never exceeds ${\approx}\SI{38}{\percent}$ utilisation and the runs are
launch-bound.

The same reading covers the other two algorithms. \emph{Per data pass},
\CaFoRand{} is \textbf{cheaper} than \BP{} (\numrange{0.76}{0.90}$\times$
Wh/epoch) and \FF{} is only \numrange{1.48}{1.76}$\times$ more expensive; the
end-to-end gaps are made almost entirely of extra epochs (\FF{} needs
\numrange{2.66}{3.82}$\times$ more, \CaFo{} \numrange{3.0}{7.8}$\times$ more).
Library maturity is not what separates these methods from \BP{}; convergence
rate is.

We note that this comment also caused us to \emph{understate} the previously
published \MF{} advantage in one respect: the submitted implementation carried
the $O(L^2)$ recomputation, so \MF{}'s cost in the accepted version was
inflated by an amount we can now bound --- and which turns out to be
negligible on this hardware.

\where{\S4.4, \emph{Is this simply an unoptimised implementation?}; \S5,
third paragraph; \S5.1, limitation~3.}

% ============================================================
\section*{Reviewer 3}
% ============================================================

\subsection*{R3a --- ``NVML \dots\ is never formally defined or explained.''}

Fixed. NVML is now expanded to the NVIDIA Management Library at first use in
the abstract and again at first use in the body, with a statement of what it
does and why it was chosen: it reports physical device state --- instantaneous
board power, SM clock, utilisation, device memory --- from the driver itself,
so an efficiency claim rests on a measured quantity rather than on an
analytical proxy such as a FLOP count.

We went further than the comment asked and \textbf{named the API call behind
every reported quantity}, because they are not interchangeable:
\texttt{nvmlDeviceGetPowerUsage} (integrated trapezoidally at \SI{5}{\hertz}
for energy), \texttt{nvmlDeviceGetUtilizationRates},
\texttt{nvmlDeviceGetClockInfo}, \texttt{nvmlDeviceGetMemoryInfo}, the host
monotonic clock for wall-clock time, \tmma{}
for per-process memory, and \texttt{psutil} for host resident set size. See
also disclosure~D5.

\where{Abstract; \S3, \emph{Evaluation metrics and instruments}; every figure
and table caption.}

\subsection*{R3b --- ``The results for CaFo in Section~5.2 are presented purely
as text \dots\ should be summarised in an equivalent table.''}

Fixed, and generalised. Rather than adding a third standalone table to a paper
already over its page limit, Table~2 now presents all three algorithms in one
structure: \FF{}, \CaFo{} (both the random-feature and the DFA-pretrained
variant) and \MF{}, each against its own tuned \BP{} baseline, with test
accuracy, training time, GPU energy and peak per-process memory as
mean~$\pm$~SD over the stated seed count. \CaFo{} therefore gets the columns
the comment asked for, and the reader gets a single point of comparison.

The \CaFo{} numbers are also no longer back-ported from an earlier cluster
software stack; all five seeds were re-measured on the current one.

\where{Table~2; \S4.2.}

\subsection*{R3c --- ``Because Mono-Forward places a cross-entropy loss at
every hidden layer, it inherently behaves like a deep supervision network
\dots\ add a standard backpropagation baseline that uses matching layer-wise
auxiliary losses. This is the only way to prove whether MF's gains stem from
its unique forward-only mechanism or simply from localised objective
constraints.''}

\textbf{Done, and the answer is a split result that we report in full.} A
\BPDS{} control was built: identical topology, identical auxiliary
cross-entropy head at every hidden layer, identical optimiser and tuning
budget --- the sole difference being that gradients flow end-to-end instead of
being detached at each stage. The table under R2c gives the outcome.

The honest reading, which is now the reading in the abstract: \emph{deep
supervision alone reproduces \MF{}'s accuracy on the MNIST-scale tasks, but
not on CIFAR, where \MF{} leads by more than \num{5}~pp.} Localised objectives
are a sufficient explanation on easy tasks; detachment is doing real work once
the task is hard enough. We have deliberately not generalised in either
direction, and we flag the mechanism (a regularisation hypothesis, grounded in
the infinite-width analysis of local losses) as a hypothesis, naming the
diagnostics that would test it and that we have not run.

We are grateful for this comment in particular: it is the one that forced the
abstract to be reframed from ``locality causes the saving'' to an attributed
decomposition.

\where{\S4.4, \emph{Does locality buy anything deep supervision does not?};
Figure~3(a), which carries all sixteen ladder contrasts with their intervals;
\S5, \emph{Why \MF{} is accurate}; abstract.}

% ============================================================
\section*{Reviewer 4}
% ============================================================

\subsection*{R4a --- ``The key term `hardware-validated comparison' is unclear
and should be explained in more detail.''}

Fixed. \S3 now opens with a definition before anything else is said: a
hardware-validated comparison is one in which \emph{every efficiency claim is
backed by a physical quantity read from the GPU itself while the workload runs
--- energy, elapsed time and memory occupancy --- rather than by an analytical
proxy such as a FLOP or parameter count.}

The operational content of that definition is supplied immediately afterwards
in the instruments paragraph (see R3a), and the specific platform is now
stated with the driver version, because the driver governs NVML power-reading
semantics: NVIDIA A100-SXM4-40GB on AMD EPYC 7742 nodes, driver 595.71.05,
Python 3.10.4, PyTorch 2.4.0+cu121.

\where{\S3, opening paragraph; \S3, \emph{Experimental Setup} and
\emph{Evaluation metrics and instruments}.}

\subsection*{R4b --- ``The statistical analysis is based on only three runs,
which is insufficient and can be considered only as preliminary evidence.''}

Accepted; this is now closed by a wide margin. Seed counts per contrast are
\FF{} $n=7$, \CaFo{} $n=5$, \MF{} ladder rungs $n=20$, and the
\BP{}/\BPDS{} ladder contrasts $n=60$ (MNIST) and $n=122$ (Fashion-MNIST) ---
\numrange{5}{122} rather than three.

We also changed the inferential frame, because more seeds without a decision
rule is just a bigger table. Every accuracy comparison is now an
\textbf{equivalence test} against a margin declared before the runs
($\pm\num{0.25}$~pp), with \num{10000}-sample paired bootstrap confidence
intervals. Of \num{28} contrasts, \textbf{eleven are equivalent, sixteen are
different, and one is inconclusive}; we report the inconclusive one as
inconclusive rather than as evidence of parity. All \num{28} appear with their
intervals in Figure~3.

We additionally report the achieved resolution so a reader can judge what the
design can and cannot detect: pooled seed-level SD of test accuracy is
\num{0.284}~pp (up to \num{0.43}~pp on the 3$\times$2000 MLPs), giving a
bootstrap CI half-width near $\pm\num{0.35}$~pp at $n=5$ and
$\pm\num{0.26}$~pp at $n=7$.

\where{\S3 protocol~5; \S4.4, \emph{Intervals, not verdicts.}; Figure~3;
\S5.1, limitation~4; seed counts in every table and caption.}

% ============================================================
\section*{Reviewer 5}
% ============================================================

\subsection*{R5a --- ``The contribution is mainly evaluative rather than
algorithmic.''}

A fair characterisation of the accepted version, and we have both sharpened
what the evaluative contribution is and identified the one place where the
work is genuinely structural rather than descriptive.

\S4.4 now presents the \textbf{time--memory frontier as a consequence of
gradient locality, not as an implementation note}. \BP{}, \BPDS{} and
\MF{}-Joint \emph{cannot} cache activations at all: joint gradients make any
cached activation stale within one optimiser step. \MF{} can, because
detachment plus layer-sequential training freeze every preceding layer for the
duration of a stage. The trade-off axis --- spend memory to avoid recomputation
--- therefore \emph{exists only for the local method}. That it turns out to be
a bad trade on this hardware is the empirical part; that it is available at all
is the algorithmic part, and it is a property of the objective structure rather
than of our code.

We have also stated the scope of the contribution plainly in \S5.1 rather than
leaving it implied: the framework, the fairness protocols, the ablation ladder
and the hardware-level measurements are ours; the three algorithms are not.

\where{\S4.4, \emph{The time--memory frontier is a structural property.};
\S5, third paragraph; \S5.1, limitation~7.}

\subsection*{R5b --- ``Small datasets, shallow architectures, and the use of
flattened CIFAR inputs for MLP experiments.''}

Accepted and retained as an explicit limitation rather than softened. \S5.1
now states that results apply to \numrange{2}{4}-layer MLPs and a 3-block CNN,
that deeper CNNs, residual networks and transformers are untested, and --- the
part we think matters most --- that flattening CIFAR to 1-D vectors destroys
spatial structure and caps absolute accuracy far below CNN state of the art,
so \emph{the comparisons are internally valid but say nothing about either
algorithm in a competitive vision pipeline}. We flatten because the \MF{}
source protocol does; we say so, and we do not claim more.

\where{\S5.1, limitations~1 and~2; \S3 protocol~1.}

\subsection*{R5c --- ``Results are based on only three runs, and some of the
reported accuracy differences are modest.''}

The first half is answered under R4b. The second half is the sharper point and
we have acted on it directly: modest differences are now \emph{adjudicated}
instead of being reported and left to the reader.

Several differences we previously reported as findings do not survive that
treatment. On MNIST 3$\times$1000 with SGD the \FF{}--\BP{} accuracy gap is
$-\num{0.014}$~pp $[-\num{0.193}, +\num{0.144}]$ and on MNIST 4$\times$2000 it
is $-\num{0.064}$~pp $[-\num{0.219}, +\num{0.070}]$; both are underpowered and
are labelled as such. \CaFoDFA{} versus \BP{} on the Fashion-MNIST CNN spans
$[-\num{0.03}, +\num{0.61}]$~pp at five seeds and is reported
\textbf{inconclusive}, not as parity. Conversely, the differences that do
matter are now unambiguous: \MF{} over \BP{} by \num{3.90} and \num{5.04}~pp
on CIFAR-10 and CIFAR-100, with intervals nowhere near zero.

\where{\S4.1; \S4.2; \S4.4, \emph{Intervals, not verdicts.}; Figure~3;
\S5.1, limitation~4.}

\subsection*{R5d --- ``A stronger comparison would include additional
baselines, such as deep-supervision or local-loss backpropagation variants, to
better isolate whether the gains come from the forward-only training mechanism
or from local auxiliary objectives.''}

Done --- this and R3c are the same request and are answered by the same
experiment. Beyond the \BPDS{} control described under R3c, the ladder adds
\textbf{\MF{}-Joint} (the \MF{} architecture, including its projection
readout, trained with end-to-end gradients), which isolates the readout from
the detachment, plus two activation-caching variants. Five rungs in total, so
that no single rung conflates two changes.

The isolation the comment asked for is therefore complete: auxiliary
objectives, the projection readout, gradient detachment and the caching
strategy are each varied alone. The answer, again, is that detachment is what
costs and --- on CIFAR only --- what pays.

\where{\S4.4; Figure~2(b), which shows the per-step contribution of all five
rungs on all four configurations.}

\subsection*{R5e --- ``The energy measurements are useful, although specific to
one GPU platform.''}

Accepted, and we have strengthened the caveat rather than merely repeating it,
because our own mechanism explanation makes it sharper than the reviewer's
phrasing. Since the central diagnosis is that these workloads are
launch-latency-bound rather than compute-bound, the results are
\emph{especially} likely to move on hardware with different kernel-launch and
memory-bandwidth characteristics. \S5.1 now states that we regard the
\textbf{direction} of the energy result, not merely its magnitude, as
platform-conditional, and the conclusion says the same.

We also disclose a second instrument boundary in the same limitation: NVML
reports GPU package power only. Host CPU, DRAM and PCIe energy are not
instrumented, so any variant that moves work off the accelerator is flattered
by the measurement. This matters most for the host-cache \MF{} strategy, whose
true total cost is understated by an unknown amount; we still report it as
dominated, since it already loses on the instrumented axis.

\where{\S5.1, limitations~3 and~5; \S5, third paragraph; \S6.}

\subsection*{R5f --- ``Clarify the limitations, and strengthen the discussion
of baselines and statistical robustness.''}

All three parts are addressed above --- limitations under R5b and R5e,
baselines under R5d, statistics under R4b and R5c. One structural change is
worth noting here: \S5.1 was rewritten from scratch. The item claiming that
three runs were a limitation is gone because it is resolved; in its place are
the achieved equivalence margin and minimum detectable effect, the GPU-only
energy accounting gap, and an explicit hyperparameter-search confound (each
algorithm was tuned independently under an equal trial budget, so differences
may partly reflect better hyperparameter discovery, and the search cost --- 
substantial for \FF{} --- is excluded from all reported figures).

\where{\S5.1 in full.}

% ============================================================
\section*{Voluntary disclosures: defects found by our own audit}
\label{sec:disclosures}
% ============================================================

\noindent
No reviewer raised any of the following. We audited our own configuration
files, logs and figures while preparing the revision, and found five defects
in the accepted version. The repository is public and cited in the paper, so
all five are discoverable by any reader after publication. We would rather
state them.

\subsection*{D1 --- Figure~3 subfloat labels were swapped (factual error in the
accepted PDF)}

In the \CaFo{} convergence float of the accepted version, the subcaption
``Validation Accuracy'' pointed at the image plotting mean validation
\emph{loss}, and ``Validation Loss'' pointed at the accuracy image. A reader
following the captions would have read both panels backwards.

The float has been removed entirely (it was one of the illegible dashboard
exports of R1b) and \CaFo{}'s results are now in Table~2, generated from the
measurement archive. Every remaining caption was checked against the image it
labels.

\subsection*{D2 --- The stopping rule was not consistent across algorithms}

\textbf{This is the single largest correction in the revision.} Protocol~4 of
the accepted version claimed ``consistent early stopping''. The configuration
files say otherwise: \BP{} stopped on validation \emph{loss}; \FF{} on
validation \emph{accuracy} with a \num{0.01} minimum improvement; \CaFo{} and
\MF{} on per-predictor and per-layer \emph{local} losses, with
per-configuration patiences and epoch caps. Worse, those caps apply per
\emph{stage}, not per run, so a nominal cap of \num{15} epochs on a four-stage
\MF{} network permits \num{60} epochs of training against \BP{}'s \num{100},
while \BP{}'s patience was in places as low as \num{5}.

Training time and energy are integrals over training duration. An inconsistent
stopping rule therefore does not perturb an efficiency comparison --- it
determines it, and it is why the \num{41}\% energy saving reverses sign once
the rule is harmonised. Table~1 gives the submitted settings against the
harmonised ones, transcribed from the configuration files, which are preserved
in the repository.

\subsection*{D3 --- Optuna pruning was asymmetric, in our own favour}

\BP{}'s Optuna studies used a median pruner; \FF{}, \CaFo{} and \MF{} ran
unpruned. \BP{}'s nominal \num{50} trials therefore included truncated ones,
and its effective search was the weaker of the two. We stress the direction:
\textbf{the asymmetry disfavoured \BP{}, i.e.\ it worked in the direction of
the accepted paper's conclusion.} Pruning is now disabled for every algorithm,
paid for by a uniform \num{20}-epoch trial cap.

\subsection*{D4 --- Three configurations were not independently searched}

The Optuna sampler was seeded from the global experiment seed, so two studies
over the same space returned byte-identical hyperparameters. The \BP{} MNIST
4$\times$2000 baseline shared its values with \BP{} MNIST 3$\times$1000, and
both \CaFo{} CIFAR-100 variants with their CIFAR-10 counterparts --- three
configurations whose hyperparameters were in effect transplanted rather than
searched. The sampler seed is now derived from the study name; all three were
re-searched, and the replacements differ by \num{2.9}$\times$,
\num{2.2}$\times$ and \num{15.3}$\times$ in at least one value. A regression
test now fails the build if any two configurations share a tuned value.

\subsection*{D5 --- Instrument misattribution, and two claims withdrawn}

Several efficiency figures in the accepted version were read off
Weights~\&~Biases dashboard exports rather than from the measurement archive,
and the captions did not name the instrument behind each quantity. Two
substantive consequences:

\begin{itemize}[nosep,leftmargin=1.4em]
  \item \textbf{Memory was reported from the wrong counter.} The
        ${\approx}\num{5}\%$ \MF{} memory saving came from the device-wide NVML
        counter, in which a fixed ${\approx}\SI{800}{\mebibyte}$ CUDA context
        --- which no algorithm controls --- dilutes every difference. Measured
        per process with \tmma{}, the
        saving is \num{27}\% on CIFAR-10 and \num{25}\% on CIFAR-100. The
        correction runs in our favour, which is not a reason to leave it
        unstated.
  \item \textbf{Two claims are withdrawn.} The accepted version attributed
        \FF{}'s slowdown in part to ``volatile GPU clock speeds \dots\
        indicating inefficient hardware saturation''. Both \BP{} and \FF{}
        hold the SM clock near \SI{1280}{\mega\hertz}; throttling is ruled
        out, not demonstrated, and the claim is gone. Likewise \MF{}'s
        ``reduced thermal profile'' --- a \SI{3}{\celsius} spread on a
        \SIrange{23}{30}{\celsius} axis --- was never evidence, and that panel
        has been removed.
\end{itemize}

\noindent
Two further corrections of the same kind: the accepted version's \FF{}
slowdown of ``up to \num{13}$\times$'' is \num{7.41}$\times$ when re-measured
under the harmonised rule, and the statement that \MF{} ``converges to a lower
final validation loss'' has been deleted because \textbf{no such series was
ever recorded} --- \MF{} logs only its layer-local objective, which is not the
network's validation loss.

Every efficiency figure in the revision is now traced to a named instrument in
\S3 and regenerated from the committed measurement archive by a committed
script.

% ============================================================
\section*{Where the revision contradicts the accepted version}
% ============================================================

For the record, and so that nothing here has to be inferred from the text:

\begin{enumerate}[nosep,leftmargin=1.6em]
  \item \textbf{Energy and time.} Accepted: \MF{} uses \num{41}\% less energy
        and trains \num{34}\% faster than \BP{}. Revised: under a harmonised
        stopping rule \MF{} is \num{4.0}$\times$ slower and uses
        \num{3.2}$\times$ the energy. The published figures reproduce under
        the published protocol; the sign reverses under a fair one.
  \item \textbf{Memory.} Accepted: ${\approx}\num{5}\%$ saving. Revised:
        \num{27}\% saving, measured per process.
  \item \textbf{\FF{} slowdown.} Accepted: up to \num{13}$\times$. Revised:
        \num{7.41}$\times$.
  \item \textbf{\FF{} clock volatility and \MF{} thermal profile.} Withdrawn.
  \item \textbf{\MF{} validation loss curve.} Withdrawn; the series does not
        exist.
  \item \textbf{\MF{} accuracy edge on CIFAR-10 under the legacy protocol.}
        Accepted: $+\num{1.21}$~pp. Re-measured at three seeds:
        $-\num{0.07}$~pp $[-\num{0.27}, +\num{0.24}]$ --- the three-seed design
        could not resolve it. Under the harmonised protocol at twenty seeds
        the edge is real and larger ($+\num{3.90}$~pp).
\end{enumerate}

\noindent
The accuracy result and the memory result stand, and the memory result is
stronger than we claimed. The efficiency result does not stand, and we say so
in the abstract, in \S4.3 and in the conclusion.

\bigskip
\noindent
We thank the reviewers again. The paper is a materially different and, we
believe, a more useful one than the version they read.

\end{document}
